\documentclass[11pt]{article}

\usepackage[margin=1in]{geometry}
\usepackage{amsmath}
\usepackage{amssymb}
\usepackage{booktabs}
\usepackage{hyperref}
\usepackage{enumitem}
\usepackage[final]{microtype}
\sloppy

\title{Round 2 NW-Sweep: Evaluating the Design Hypotheses Against Real Data}
\author{FastMSPEC project report}
\date{2026-09-04}

\begin{document}
\maketitle

\begin{abstract}
Round 2 (`docs/notebook5\_revamp\_progress.md`'s "Round 2 design" section) ran FastMspec at 5
candidate bandwidths per pair for a 60-pair, distance-stratified subset (300 work units total),
built specifically to test four hypotheses that motivated the sweep's design in the first place.
This report states each hypothesis plainly, evaluates it against the real 300-point dataset, and
gives an explicit verdict for each -- including one result (H2) that did \emph{not} come out the
way a naive reading of the theory doc would predict, reported as found rather than smoothed over.
Standalone reference, independent of the conversation that produced it.
\end{abstract}

\tableofcontents

\section{The Four Hypotheses, Stated for the Record}
\label{sec:hypotheses}

Recalled from the session that designed Round 2; stated here explicitly so any misremembering
can be corrected before the analysis below is trusted.

\begin{description}[leftmargin=1.5em]
  \item[H1 --- Normalized bandwidth is the right variable.] A pair's own $NW_{\rm high}(r) \approx
    Nc_{\min}/(4r)$ (the resolution-bandwidth ceiling, computed per pair from its distance and
    local phase velocity) is the correct quantity to normalize candidate bandwidths against --
    i.e., convergence/quality should behave as a consistent function of $NW/NW_{\rm high}(r)$
    across pairs of very different absolute distance, not depend on absolute $NW$ in a way that
    differs pair to pair.
  \item[H2 --- An interior MSE optimum exists.] Per the bias--variance framework in
    \texttt{stage5\_bandwidth\_theory.tex} (following Haley \& Anitescu 2017), quality should
    \emph{peak} at some bandwidth strictly between "as narrow as possible" and the resolution
    ceiling -- not simply keep improving all the way to the ceiling. Haley \& Anitescu's own
    Lorentzian example landed at $NW=14.5$ against a much larger ceiling, motivating the
    expectation of an interior peak here too.
  \item[H3 --- An existence threshold produces a cliff for far pairs.] For sufficiently large
    distance $r$ (hence small $NW_{\rm high}(r)$), no candidate bandwidth in a reasonable range
    should produce reliable convergence at all -- a qualitatively different, complete-failure
    pattern for far pairs, distinct from the gradual bias--variance degradation expected for near
    pairs.
  \item[H4 --- A hard numerical floor exists below $NW\approx 3$.] Discovered mid-round, not
    originally hypothesized, but directly testable once found: \texttt{FastMultitaper}'s taper
    count $K$ hits exactly zero below $NW\approx 3$ (this project's \texttt{cutoff}), a
    structurally different failure mode from H2/H3's statistical ones, expected to concentrate at
    low fractions and, because $NW_{\rm high}(r)$ shrinks with $r$, disproportionately on far
    pairs at any \emph{fixed} fraction.
\end{description}

\section{Data}

300 work units: 60 pairs (15 per distance quartile, using Round 1's own quartile boundaries) x 5
bandwidth fractions each, $NW/NW_{\rm high}(r) \in \{1.5, 1.0, 0.7, 0.45, 0.25\}$. Breakdown:
\textbf{60/300 converged (20\%)}, \textbf{65/300 (21.7\%) hit the H4 floor} ($K=0$, a clear,
diagnosable error after the same-day fix -- see \texttt{docs/notebook5\_revamp\_progress.md}'s
2026-09-04 log entry), \textbf{175/300} are genuine convergence/non-convergence outcomes at
numerically valid bandwidths. All numbers below are computed directly from
\texttt{data/results/dispcurve\_quality/round2\_sweep/} (300 files) and \texttt{round2\_subset.csv}.

\section{H1: Is fraction-of-$NW_{\rm high}$ the right normalization?}

\begin{table}[h]
\centering
\begin{tabular}{lccccc}
\toprule
Quartile & $NW/NW_{\rm high}=1.5$ & $1.0$ & $0.7$ & $0.45$ & $0.25$ \\
\midrule
1 (nearest) & 12/15 (80\%) & 10/15 (67\%) & 9/15 (60\%) & 7/15 (47\%) & 3/15 (20\%) \\
2           & 8/15 (53\%)  & 4/15 (27\%)  & 2/15 (13\%) & 0/15 (0\%)  & 0/15 (0\%) \\
3           & 2/15 (13\%)  & 3/15 (20\%)  & 0/15 (0\%)  & 0/15 (0\%)  & 0/15 (0\%) \\
4 (farthest)& 0/15 (0\%)   & 0/15 (0\%)   & 0/15 (0\%)  & 0/15 (0\%)  & 0/15 (0\%) \\
\bottomrule
\end{tabular}
\caption{Convergence count/rate by quartile x fraction (raw, $K=0$ rows included as
non-converged).}
\label{tab:h1}
\end{table}

\textbf{Verdict: partially supported, with an important qualification.} Quartiles 1--3 all show
the same \emph{direction} of dependence on fraction (convergence declines as fraction shrinks),
consistent with $NW/NW_{\rm high}(r)$ being a meaningful normalized variable rather than an
artifact of absolute distance. But quartile 4 doesn't participate in this pattern at all -- it is
uniformly zero across every fraction tested, including $1.5\times$ its own ceiling. Normalizing
by $NW_{\rm high}(r)$ makes quartiles 1--3 comparable to each other; it does not rescue quartile
4, whose true operating point (if one exists at all for this window length) appears to sit
entirely outside the tested range. This is exactly what H3 predicts, and is evaluated properly
there rather than being treated as a failure of H1.

\section{H2: Is there an interior optimum bandwidth?}

\begin{table}[h]
\centering
\begin{tabular}{lccccc}
\toprule
& $1.5$ & $1.0$ & $0.7$ & $0.45$ & $0.25$ \\
\midrule
Convergence rate, raw (n=60 each) & 37\% & 28\% & 18\% & 12\% & 5\% \\
Convergence rate, excl. $K=0$ rows & 37\% & 28\% & 21\% & 17\% & 14\% \\
Mean \texttt{bad\_quality\_fraction} (converged only) & 0.634 & 0.651 & 0.663 & 0.666 & 0.800 \\
Mean \texttt{freq\_coverage\_fraction} (converged only) & 0.540 & 0.486 & 0.476 & 0.358 & 0.366 \\
\bottomrule
\end{tabular}
\caption{Convergence rate and quality scalars vs. fraction, aggregated across all quartiles.}
\label{tab:h2}
\end{table}

\textbf{Verdict: not supported within the tested range $[0.25, 1.5]\times NW_{\rm high}$.} Every
measure -- raw convergence, $K=0$-excluded convergence, and both quality scalars among converged
units -- moves monotonically in the same direction: wider is better, all the way out to
$1.5\times$ the nominal resolution ceiling, with no sign of turning over. This is not the pattern
H2 predicted. Two honest readings, not adjudicated here:
\begin{enumerate}[nosep]
  \item The true interior optimum lies \emph{above} $1.5\times NW_{\rm high}$ -- i.e., the
    resolution ceiling itself is more conservative than the point where bias actually starts
    dominating, and the sweep simply didn't reach far enough to see the turnover.
  \item In this regime, variance (favoring wider bandwidth) dominates the smooth bias term badly
    enough, or other project-specific effects (picker mechanics, real non-stationarity per
    \texttt{stage5\_bandwidth\_theory.tex}'s Section 5.2) swamp it enough, that no interior
    optimum is visible at this sample size even if one exists in principle.
\end{enumerate}
Distinguishing these needs a wider sweep (fractions above 1.5) before H2 can be considered settled
either way -- reported as an open question, not resolved by extrapolation.

\section{H3: Does a real existence-threshold cliff appear for far pairs?}

\textbf{Verdict: strongly supported, and more starkly than the qualitative prediction.} Quartile
4's 15 pairs (individually, not just in aggregate) show \emph{zero} converged units across all 75
attempts (5 fractions x 15 pairs), including at $1.5\times$ their own already-small $NW_{\rm
high}(r)$ (range 2.30--5.08, all pairs individually listed in the underlying data). This is not a
gradual decline -- it's a complete absence of convergence across the entire tested bandwidth
range for every single far pair, exactly the qualitative "no valid window exists" signature H3
predicted, and clearer than a marginal cliff would have been.

\section{H4: The $K=0$ numerical floor}

\begin{table}[h]
\centering
\begin{tabular}{lcccc}
\toprule
Quartile 1 & Quartile 2 & Quartile 3 & Quartile 4 \\
\midrule
1/75 (1\%) & 9/75 (12\%) & 17/75 (23\%) & 38/75 (51\%) \\
\bottomrule
\end{tabular}
\caption{$K=0$-floor rate by quartile (out of all 75 sweep points per quartile).}
\label{tab:h4}
\end{table}

\textbf{Verdict: confirmed, with a clean, mechanistically-explained gradient.} The $K=0$ floor
rate rises monotonically from 1\% (nearest quartile) to 51\% (farthest) -- exactly what the
mechanism predicts: since $NW_{\rm high}(r)$ shrinks with distance, the \emph{same} fraction
multiplier produces a smaller absolute $NW$ for farther pairs, crossing the $NW\approx 3$ floor at
a higher fraction. By quartile 4, over half of all sweep points are not merely statistically
unreliable but numerically undefined. This is a hard, structural limit, categorically different
from H2/H3's statistical ones -- restated from \texttt{docs/notebook5\_revamp\_progress.md}: it
must not be conflated with the soft, bias-variance-driven lower bound Stage 5 is deriving.

\section{Synthesis and Implications for Stage 5}

\begin{itemize}[nosep]
  \item \textbf{H1 and H3 together give a clean, load-bearing result}: normalized bandwidth
    ($NW/NW_{\rm high}(r)$) is the right variable for pairs where a viable window exists at all
    (quartiles 1--3), and the existence-threshold theory correctly predicts \emph{which} pairs
    don't have one (quartile 4, confirmed with zero exceptions). This validates the core
    machinery \texttt{round2\_prep.py} was built on, not just the sweep's raw convergence numbers.
  \item \textbf{H2's non-result is the most actionable finding for Stage 5's next step}: the
    bias--variance MSE framework predicts an interior optimum, and this sweep didn't find one --
    not because it was falsified, but because the tested range may not have reached it. Before
    trusting \texttt{stage5\_bandwidth\_theory.tex}'s Eq.~16 recipe as a complete answer, extend
    the sweep's upper fraction (e.g. to $2.0$--$3.0\times NW_{\rm high}$) on a subset of quartile-1
    pairs specifically, where a real signal is most likely to be visible against the noise floor.
  \item \textbf{H4 is now a known, permanent constraint on any future sweep design}: any bandwidth
    search must stay above $NW\approx 3$ or expect the same clean failure mode -- already fixed to
    fail informatively rather than silently, but still worth excluding from statistical fits by
    construction (e.g. floor the fraction grid's minimum well above wherever $NW\approx 3$ falls
    for the pairs being tested) rather than discovered again.
\end{itemize}

\section{Caveats}

\begin{itemize}[nosep]
  \item Sample sizes are modest -- 15 pairs per quartile, and after excluding the $K=0$ floor,
    as few as 22 valid units at the lowest fraction. The quartile-level patterns (Tables
    \ref{tab:h1}, \ref{tab:h4}) are stark enough to trust qualitatively; finer claims (e.g. exact
    convergence-rate percentages) carry real sampling uncertainty at this $n$.
  \item Quality scalars (Table \ref{tab:h2}, bottom two rows) are averaged only over
    \emph{converged} units at each fraction -- as convergence rate drops, this average is
    increasingly selected on a shrinking, potentially biased subset (as few as 3 units at
    fraction $0.25$), not a random sample of "quality regardless of outcome."
  \item This report covers FastMspec only, the technique the sweep was designed around;
    single-taper's own Round 1 result (0/380 converged, a real mechanistic result) is separate
    context, not re-evaluated here.
\end{itemize}

\end{document}
